\section{実験装置}
本実験では以下の装置を用いた。
\begin{itemize}
  \item \textbf{ファンクションジェネレータ(FC)}(メーカー: NF, 型式: WF1973, SER: 9219781)\\
  回路へ入力電圧を加えるために使用した。
  \item \textbf{ディジタルマルチメータ(DMM)}(メーカー: KEYSIGHT, 型式: 34450A, SER: MY57410006)\\
  回路の抵抗に加わる電圧を測定するために使用。
  \item \textbf{バッファーアンプ}\\
  入力されたシグナルを伝送に適したシグナルに変換するアンプ。
  \item \textbf{ブレッドボード}\\
  回路素子を取り付け回路を作る基盤。
  \item \textbf{シリコンダイオード}(1S1588)\\
  静特性の測定対象。実験1~4で使用。 
  \item \textbf{ゲルマニウムダイオード}(1N60)\\
  静特性の測定対象。実験4で使用。
  \item \textbf{ツェナーダイオード}(RD5.6E)\\
  静特性の測定対象。実験4で使用。
  \item \textbf{抵抗}\\
  ダイオードの特性を調べるためダイオードと直列に接続。
  \item \textbf{PC}\\
  LabVIEWを用いて自動計測のプログラムを作成し,データの取得も行った。
\end{itemize}

